% GENERATED FILE. Edit canonical lessons, then run npm run generate:books.
% canonical-source-hash: fnv1a64:399b29981051b086
% canonical-lessons: FA-C492-halat-e-tahavvo, FA-C492-estefragh-kardan, FA-C492-ghash-kardan, FA-C492-zaf-kardan, FA-C492-araq-kardan

\chapter{Feeling Unwell}
\label{ch:fa-feeling-unwell}
\hlchaptermodality{492}

\begin{chapteropening}
\textbf{By the end of this chapter:} \emph{I can say that I feel sick or faint.}

Complete the last lesson of chapter 492: I can say that I feel sick or faint.
\end{chapteropening}

\section[\texorpdfstring{\fa{حالت} \fa{تهوع} (hâlat-e tahavvo')}{حالت تهوع (hâlat-e tahavvo')}]{\fa{حالت} \fa{تهوع} (hâlat-e tahavvo') — nausea}

\label{lesson:FA-C492-halat-e-tahavvo}



\begin{quote}
Before the new one: say the Persian for an itch, then the Persian for to hiccup.
\end{quote}

\subsection*{You'll want to know: \fa{حالت} \fa{تهوع}}
\textbf{\fa{حالت} \fa{تهوع}} (\emph{hâlat-e tahavvo'}) — \textquotedblleft{}nausea\textquotedblright{}.

\begin{quote}
\fa{در} \fa{ماشین} \fa{حالت} \fa{تهوع} \fa{میگیرم}.
\end{quote}

\begin{quote}
\emph{dar mâshin hâlat-e tahavvo' migiram.}
\end{quote}

\begin{quote}
I feel sick in the car.
\end{quote}

\emph{hâlat} 'state' + \emph{tahavvo'} 'retching'. \emph{hâlat-e tahavvo' dâshtan} — to feel sick.

\begin{cousinweb}
\emph{hâlat-e tahavvo'} is built from two Arabic loans: \emph{hâlat}, from Arabic \emph{ḥāla}, 'state' (EtymDB), and \emph{tahavvo'}, usually traced to Arabic \emph{tahawwuʿ}, 'retching'.
\end{cousinweb}

\begin{grammarlens}[title={What you've built}]
One more word for telling, explaining and describing at B1.
\end{grammarlens}

\subsection*{Your turn}
\begin{itemize}
\raggedright
  \item \emph{Say it:} \emph{hâlat-e tahavvo'}
  \item \emph{Say it:} \emph{hâlat-e tahavvo'}, once more
  \item \emph{Say it:} a sentence of your own with \emph{hâlat-e tahavvo'}
  \item \emph{Recall:} say the Persian for an itch, then the Persian for to hiccup, then say \emph{hâlat-e tahavvo'} again
\end{itemize}

\begin{tcolorbox}[breakable,colback=teal!4,colframe=teal!35!black,title={Before you move on}]
What does \fa{حالت} \fa{تهوع} mean? (\textquotedblleft{}Nausea\textquotedblright{}.) Say it once more.
\end{tcolorbox}
\section[\texorpdfstring{\fa{استفراغ} \fa{کردن} (estefrâgh kardan)}{استفراغ کردن (estefrâgh kardan)}]{\fa{استفراغ} \fa{کردن} (estefrâgh kardan) — to vomit, to be sick}

\label{lesson:FA-C492-estefragh-kardan}



\begin{quote}
Before the new one: say the Persian for to hiccup, then the Persian for nausea.
\end{quote}

\subsection*{You'll want to know: \fa{استفراغ} \fa{کردن}}
\textbf{\fa{استفراغ} \fa{کردن}} (\emph{estefrâgh kardan}) — \textquotedblleft{}to vomit, to be sick\textquotedblright{}.

\begin{quote}
\fa{بچه} \fa{در} \fa{اتوبوس} \fa{استفراغ} \fa{کرد}.
\end{quote}

\begin{quote}
\emph{bachche dar otobus estefrâgh kard.}
\end{quote}

\begin{quote}
The child was sick on the bus.
\end{quote}

\emph{estefrâgh} 'vomiting' + \emph{kardan}. More medical than \emph{bâlâ âvardan}.

\begin{cousinweb}
\emph{estefrâgh} is a loan, usually traced to Arabic \emph{istifrāġ}, 'emptying out'. EtymDB has no entry for it. \emph{kardan}, 'to do', is inherited Persian.
\end{cousinweb}

\begin{grammarlens}[title={What you've built}]
One more word for telling, explaining and describing at B1.
\end{grammarlens}

\subsection*{Your turn}
\begin{itemize}
\raggedright
  \item \emph{Say it:} \emph{estefrâgh kardan}
  \item \emph{Say it:} \emph{estefrâgh kardan}, once more
  \item \emph{Say it:} a sentence of your own with \emph{estefrâgh kardan}
  \item \emph{Recall:} say the Persian for to hiccup, then the Persian for nausea, then say \emph{estefrâgh kardan} again
\end{itemize}

\begin{tcolorbox}[breakable,colback=teal!4,colframe=teal!35!black,title={Before you move on}]
What does \fa{استفراغ} \fa{کردن} mean? (\textquotedblleft{}To vomit, to be sick\textquotedblright{}.) Say it once more.
\end{tcolorbox}
\section[\texorpdfstring{\fa{غش} \fa{کردن} (ghash kardan)}{غش کردن (ghash kardan)}]{\fa{غش} \fa{کردن} (ghash kardan) — to faint}

\label{lesson:FA-C492-ghash-kardan}



\begin{quote}
Before the new one: say the Persian for nausea, then the Persian for to vomit.
\end{quote}

\subsection*{You'll want to know: \fa{غش} \fa{کردن}}
\textbf{\fa{غش} \fa{کردن}} (\emph{ghash kardan}) — \textquotedblleft{}to faint\textquotedblright{}.

\begin{quote}
\fa{او} \fa{در} \fa{گرما} \fa{غش} \fa{کرد}.
\end{quote}

\begin{quote}
\emph{u dar garmâ ghash kard.}
\end{quote}

\begin{quote}
She fainted in the heat.
\end{quote}

\emph{ghash} 'a faint' + \emph{kardan}. Colloquially also 'to laugh till you drop': \emph{az khande ghash kard}.

\begin{cousinweb}
\emph{ghash} is a loan, usually traced to Arabic \emph{ġašy}, 'a fainting fit'. EtymDB has no entry for it. \emph{kardan}, 'to do', is inherited Persian.
\end{cousinweb}

\begin{grammarlens}[title={What you've built}]
One more word for telling, explaining and describing at B1.
\end{grammarlens}

\subsection*{Your turn}
\begin{itemize}
\raggedright
  \item \emph{Say it:} \emph{ghash kardan}
  \item \emph{Say it:} \emph{ghash kardan}, once more
  \item \emph{Say it:} a sentence of your own with \emph{ghash kardan}
  \item \emph{Recall:} say the Persian for nausea, then the Persian for to vomit, then say \emph{ghash kardan} again
\end{itemize}

\begin{tcolorbox}[breakable,colback=teal!4,colframe=teal!35!black,title={Before you move on}]
What does \fa{غش} \fa{کردن} mean? (\textquotedblleft{}To faint\textquotedblright{}.) Say it once more.
\end{tcolorbox}
\section[\texorpdfstring{\fa{ضعف} \fa{کردن} (za'f kardan)}{ضعف کردن (za'f kardan)}]{\fa{ضعف} \fa{کردن} (za'f kardan) — to feel faint (from hunger)}

\label{lesson:FA-C492-zaf-kardan}



\begin{quote}
Before the new one: say the Persian for to vomit, then the Persian for to faint.
\end{quote}

\subsection*{You'll want to know: \fa{ضعف} \fa{کردن}}
\textbf{\fa{ضعف} \fa{کردن}} (\emph{za'f kardan}) — \textquotedblleft{}to feel faint (from hunger)\textquotedblright{}.

\begin{quote}
\fa{صبحانه} \fa{نخوردم} \fa{و} \fa{ضعف} \fa{کردم}.
\end{quote}

\begin{quote}
\emph{sobhâne nakhordam va za'f kardam.}
\end{quote}

\begin{quote}
I skipped breakfast and felt faint.
\end{quote}

\emph{za'f} 'weakness' + \emph{kardan}. The adjective is \emph{za'if} 'weak'.

\begin{cousinweb}
\emph{za'f} is a loan, usually traced to Arabic \emph{ḍaʿf}, 'weakness'. EtymDB has no entry for it. \emph{kardan}, 'to do', is inherited Persian.
\end{cousinweb}

\begin{grammarlens}[title={What you've built}]
One more word for telling, explaining and describing at B1.
\end{grammarlens}

\subsection*{Your turn}
\begin{itemize}
\raggedright
  \item \emph{Say it:} \emph{za'f kardan}
  \item \emph{Say it:} \emph{za'f kardan}, once more
  \item \emph{Say it:} a sentence of your own with \emph{za'f kardan}
  \item \emph{Recall:} say the Persian for to vomit, then the Persian for to faint, then say \emph{za'f kardan} again
\end{itemize}

\begin{tcolorbox}[breakable,colback=teal!4,colframe=teal!35!black,title={Before you move on}]
What does \fa{ضعف} \fa{کردن} mean? (\textquotedblleft{}To feel faint (from hunger)\textquotedblright{}.) Say it once more.
\end{tcolorbox}
\section[\texorpdfstring{\fa{عرق} \fa{کردن} ('araq kardan)}{عرق کردن ('araq kardan)}]{\fa{عرق} \fa{کردن} ('araq kardan) — to sweat}

\label{lesson:FA-C492-araq-kardan}



\begin{quote}
Before the new one: say the Persian for to faint, then the Persian for to feel faint.
\end{quote}

\subsection*{You'll want to know: \fa{عرق} \fa{کردن}}
\textbf{\fa{عرق} \fa{کردن}} (\emph{'araq kardan}) — \textquotedblleft{}to sweat\textquotedblright{}.

\begin{quote}
\fa{در} \fa{باشگاه} \fa{خیلی} \fa{عرق} \fa{کردم}.
\end{quote}

\begin{quote}
\emph{dar bâshgâh kheyli 'araq kardam.}
\end{quote}

\begin{quote}
I sweated a lot at the gym.
\end{quote}

\emph{'araq} 'sweat' + \emph{kardan}. \emph{'araq} also names a distilled herbal water, such as \emph{'araq-e na'nâ} 'mint water'.

\begin{cousinweb}
\emph{'araq} comes from Arabic \emph{ʿaraq}, 'sweat' (EtymDB). \emph{kardan}, 'to do', is inherited Persian.
\end{cousinweb}

\begin{grammarlens}[title={What you've built}]
One more word for telling, explaining and describing at B1.
\end{grammarlens}

\subsection*{Your turn}
\begin{itemize}
\raggedright
  \item \emph{Say it:} \emph{'araq kardan}
  \item \emph{Say it:} \emph{'araq kardan}, once more
  \item \emph{Say it:} a sentence of your own with \emph{'araq kardan}
  \item \emph{Recall:} say the Persian for to faint, then the Persian for to feel faint, then say \emph{'araq kardan} again
\end{itemize}

\begin{tcolorbox}[breakable,colback=teal!4,colframe=teal!35!black,title={Before you move on}]
What does \fa{عرق} \fa{کردن} mean? (\textquotedblleft{}To sweat\textquotedblright{}.) Say it once more.
\end{tcolorbox}
